%% kapa: exam=1 round=32 year=2021 period=1 form=A subject=economics subject_ko=경제학원론 seq=5 number=45
%% source: Q-Net 감정평가사 기출문제 (cst00309) artlSeq=5208397 file=extracted/1008/5208397/A형/2021년도 제32회 감정평가사 1차 1교시 A형.pdf
%% conversion: kordoc 4.24.51 -> tools/gen1.py (auto); review=proofread-sonnet; flags=none
\begin{question}{45}
수요곡선은 $P=10$, 공급곡선은 $Q_s=P$이다. 정부가 한 단위당 2원의 물품세를 소비자에게 부과한 결과로 옳은 것은? (단, $Q_s$는 공급량, $P$는 가격이다.)
\begin{choices}
  \item 소비자 대신 생산자에게 물품세를 부과하면 결과는 달라진다.
  \item 소비자잉여는 감소하였다.
  \item 생산자잉여의 감소분은 24원이다.
  \item 자중손실(deadweight loss)은 2원이다.
  \item 조세수입은 20원 증가하였다.
\end{choices}
\end{question}
\answer{4}
